\section{Conclusão e trabalhos futuros}

Reconstruímos do ChEMBL\_37 uma base de 27.714 moléculas ativas contra EGFR, BACE1, DRD2 e HDAC1 e treinamos uma árvore de decisão para prever o alvo. Com divisão agrupada, os cinco descritores físico-químicos dão acurácia de 0,771; acrescentar 12 contagens de caracteres do SMILES a leva a 0,894~$\pm$~0,004, e as contagens sozinhas já chegam a 0,876. A divisão aleatória superestima o resultado em 1,3 a 4,1 pontos porque cerca de um quarto das moléculas de teste tem gêmea de descritores no treino. TPSA e o número de átomos aromáticos são os atributos mais importantes, e o par HDAC1--EGFR concentra 28,7\% dos erros. A Random Forest chega a 0,952, cerca de 6 pontos acima da árvore, que em troca oferece regras legíveis.

Quatro caminhos seguem do que foi visto. Fingerprints estruturais, calculados com uma biblioteca de química, enxergariam subestruturas como o hidroxamato, que as contagens não veem. Uma divisão por esqueleto molecular daria uma estimativa menos otimista para moléculas novas. A curva de aprendizado ainda sobe, o que sugere buscar mais dados no ChEMBL, com outros alvos ou um limiar de atividade mais frouxo. Por fim, podar a árvore ou usar um conjunto de árvores daria probabilidades calibradas, e incluir moléculas inativas faria o modelo servir também à triagem.
